\documentclass[11pt]{article}

\usepackage[T1]{fontenc}
\usepackage[utf8]{inputenc}
\usepackage{lmodern}
\usepackage[margin=1.06in]{geometry}
\usepackage{amsmath,amssymb,amsthm,mathtools}
\usepackage{enumitem}
\usepackage{microtype}
\usepackage{xcolor}
\usepackage{hyperref}

\hypersetup{
  colorlinks=true,
  linkcolor=blue!45!black,
  citecolor=green!35!black,
  urlcolor=blue!55!black,
  pdftitle={Supplementary Notes for the Affine-Line Exponential-Period Theorem},
  pdfauthor={Christopher D. Long}
}

\newtheorem{proposition}{Proposition}[section]
\newtheorem{lemma}[proposition]{Lemma}
\newtheorem{remark}[proposition]{Remark}
\theoremstyle{definition}
\newtheorem{definition}[proposition]{Definition}

\newcommand{\Qbar}{\overline{\mathbb Q}}
\newcommand{\C}{\mathbb C}
\newcommand{\Aone}{\mathbb A^1}
\newcommand{\Pcal}{\mathcal P}
\newcommand{\Bcal}{\mathcal B}
\newcommand{\Ocal}{\mathcal O}
\newcommand{\supp}{\operatorname{supp}}

\title{Supplementary Notes for the Affine-Line\
Exponential-Period Theorem}
\author{Christopher D. Long}
\date{August 2026}

\begin{document}
\maketitle

\begin{abstract}
These notes preserve constructions removed from the streamlined Paper~II,
\emph{A Linear Kontsevich--Zagier Theorem for Compact Exponential Periods
on the Affine Line}.  They are not needed for its shortest proof.  The most
useful retained items are an explicit boundary presentation suited to
computation, affine covariance, and an abstract spectral criterion behind
direct-sum horizontal rigidity.
\end{abstract}

\section{The explicit boundary presentation}

Let $\Delta=\sum_\alpha d_\alpha[\alpha]$ be a finite degree-zero
zero-cycle on $\Aone(\Qbar)$ with coefficients in $\Qbar$.  For
$q,r\in\Qbar[x]$, write
\[
  \mathsf P(q,r,\Delta)
\]
for a formal compact affine-line exponential period, and write
$\mathsf E(\beta)$ for the zero-dimensional exponential period $e^\beta$.

\begin{definition}[Boundary presentation]
Let $\Bcal_{\Aone}^{\mathrm{cpt}}$ be the $\Qbar$-vector space generated by
these symbols, modulo linearity in $r$ and $\Delta$ and the following
relation.  Let $Y/\Qbar$ be a smooth connected affine curve, let
$\pi\in\Ocal(Y)$, and let $\eta$ be a compact relative one-chain on
$Y(\C)$ with algebraic boundary
\[
  \partial\eta=\sum_p n_p[p].
\]
Suppose that morphisms $\rho_\nu:Y\to\Aone$, polynomials
$q_\nu,r_\nu\in\Qbar[x]$, and coefficients $c_\nu\in\Qbar$ satisfy
\[
  q_\nu\circ\rho_\nu=\pi
\]
and
\[
  \sum_\nu c_\nu\rho_\nu^*(r_\nu(x)\,dx)
  =dG+G\,d\pi
\]
for some $G\in\Ocal(Y)$.  Put
$\Delta_\nu=\partial(\rho_{\nu*}\eta)$.  Impose
\begin{equation}\label{eq:boundary-relation}
  \sum_\nu c_\nu\mathsf P(q_\nu,r_\nu,\Delta_\nu)
  =\sum_p n_pG(p)\mathsf E(\pi(p)).
\end{equation}
\end{definition}

Relation \eqref{eq:boundary-relation} simultaneously encodes change of
variables and twisted Stokes.  It gives a well-defined period map
\[
  \mathsf P(q,r,\Delta)
  \longmapsto
  \int_\Gamma r(x)e^{q(x)}\,dx,
  \qquad
  \mathsf E(\beta)\longmapsto e^\beta,
\]
where $\partial\Gamma=\Delta$.

There is a canonical surjection from this boundary presentation to the
affine-line subspace of the curve-level formal period space used in the
streamlined paper.  Indeed, \eqref{eq:boundary-relation} follows there from
functoriality, bilinearity, and twisted Stokes.

\begin{proposition}[Equivalence of the two presentations]
Assuming the main theorem of Paper~II, the boundary presentation and the
affine-line subspace of the curve-level formal period space are canonically
isomorphic.
\end{proposition}

\begin{proof}
The canonical map is surjective.  The proof of Paper~II shows directly that
a functional relation has vanishing inverse-branch densities.  On the
simultaneous Galois cover, each edge relation has the form
\eqref{eq:boundary-relation} with $G=0$; twisted reduction has the same form
with $Y=\Aone$.  Therefore the arithmetic transfer theorem proves that the
period map on the boundary presentation is injective.  Its canonical
surjection to the curve-level subspace consequently has zero kernel.
\end{proof}

The boundary presentation is longer and less standard, which is why it was
removed from the main paper.  It may nevertheless be useful for an
algorithm that stores relations entirely in terms of affine-line boundary
data and finite algebraic covers.

\section{Affine covariance}

\begin{lemma}[Affine covariance]
Let $a\in\Qbar^\times$, $b\in\Qbar$, and put
\[
  \widetilde q(y)=q(ay+b),
  \qquad
  \widetilde r(y)=a\,r(ay+b),
  \qquad
  \widetilde\Delta
  =\sum_\alpha d_\alpha
    \left[\frac{\alpha-b}{a}\right].
\]
Then
\[
  I_{q,r,\Delta}(z)
  =I_{\widetilde q,\widetilde r,\widetilde\Delta}(z)
\]
for all $z$, and the corresponding formal symbols are equal.
\end{lemma}

\begin{proof}
This is the change of variables $x=ay+b$.  Formally, it is the
functoriality relation for the automorphism $y\mapsto ay+b$ of $\Aone$.
\end{proof}

This lemma remains useful for examples or normal forms, but it should not be
used to normalize the polynomial in the proof: the unnormalized moment
system already has the required residue spectrum, and normalization only
introduces extra symbols and transformed cycles.

\section{An abstract direct-sum rigidity criterion}

The direct-sum argument depends only on spectral nonresonance, not on the
particular origin of the moment systems.

\begin{proposition}[Spectral criterion]
Let $C,D\in M_n(\C)$, with $C$ semisimple.  Assume:
\begin{enumerate}[label=\textup{(\roman*)},leftmargin=2.4em]
\item no negative integer is an eigenvalue of
$\operatorname{ad}D:X\mapsto[D,X]$;
\item for every eigenvalue $\lambda$ of $C$, no positive integer is an
eigenvalue of
\[
  X\longmapsto
  [\pi_\lambda D|_{\ker(C-\lambda I)},X].
\]
\end{enumerate}
If $P\in M_n(\C(z))$ satisfies
\[
  P'=[C+D/z,P],
\]
then $P$ is constant and commutes with $C,D$.
\end{proposition}

\begin{proof}
A pole away from zero is excluded by comparing pole orders.  At zero, the
leading Laurent coefficient of a pole of order $m$ would be a
$-m$-eigenvector of $\operatorname{ad}D$, contrary to (i).  Thus $P$ is a
polynomial.  If its degree is $d>0$, its leading coefficient commutes with
$C$.  Compressing the next coefficient equation to a $C$-eigenspace makes a
nonzero leading block a $d$-eigenvector of the commutator in (ii), a
contradiction.
\end{proof}

A convenient sufficient condition is that the eigenvalues of $D$ lie in an
interval of length less than one, and that the eigenvalues of every
compressed block lie in an interval of length less than one.  The polynomial
moment systems satisfy the stronger condition that all relevant spectra lie
in $(-1,0)$.  This abstract form may be useful for other finite maps or
other Fourier--Laplace connections.

\section{Algebraic evaluation points}

For $\xi\in\Qbar^\times$,
\[
  I_{q,r,\Delta}(\xi z)=I_{\xi q,r,\Delta}(z).
\]
Thus a finite relation involving separate algebraic arguments $\xi_j$ is
already a relation for the polynomials $\xi_jq_j$.  Carrying the $\xi_j$ as
additional theorem parameters obscures this fact and is unnecessary.

There is one qualification.  Twisted reduction of a general polynomial
differential at the numerical value $z=\xi$ uses
$R'+\xi q'R$, not $R'+q'R$.  The clean procedure is therefore to absorb
$\xi$ into $q$ first and then perform the canonical reduction for the new
polynomial.

\section{Terminology concordance}

The streamlined papers use the vocabulary of the exponential-period,
$E$-function, and polynomial-moment literature.  The following replacements
were made systematically:
\begin{center}
\begin{tabular}{ll}
Earlier draft & Streamlined terminology\\[2pt]
\hline
polynomial phase & polynomial $q$ or regular function\\
phase value & value of $q$\\
multiphase relation & relation involving several polynomials\\
phase moment & moment in $q$\\
phase density & inverse-branch density\\
common constellation & common tree representation\\
zero-chain & degree-zero zero-cycle\\
polynomial-line category & affine-line subspace\\
amplitude & polynomial differential or coefficient polynomial\\
evaluation parameter $\xi$ & absorbed into $q$\\
\end{tabular}
\end{center}

The word \emph{phase} remains appropriate in a stationary-phase, oscillatory,
or Stokes-asymptotic argument.  It is not needed as the generic name for the
polynomial in these two papers.

\section{Items deliberately not promoted to the main theorem}

The explicit boundary presentation, affine normalization, and separate
evaluation parameters are all valid.  They were removed because they add
formal overhead without strengthening the central statement.  The abstract
spectral criterion is more reusable, but including it in the main paper
would interrupt the proof's one-dimensional narrative.  These notes retain
all four ideas for later algorithmic or geometric extensions.

\end{document}
